% Values: PRECISE, FAITHFUL, ILLUMINATING, TRACTABLE, ECONOMICAL
% Evidence planned:
% - PRECISE: the syntax of programs with its restrictions; one rule for each
%   production of Program in elaboration and in evaluation, and one rule
%   combining them over Basis = StaticBasis x DynamicBasis; configurations,
%   events and every component of a configuration defined by set
%   constructions; dispatch, the start of a program and every transition
%   (update, process steps and their actions, ends of processes, incoming and
%   outgoing port values, the passing of time, platform steps) given by rules;
%   the meaning of a program as a set of traces, and conformance.
% - FAITHFUL: A32 and A82 (only the main module's main is constrained; it
%   need not be exposed; Program or a view); A27 (Dflt); A84 (port names
%   global, static); A72 and A77 (flags); A66, A67 (dispatch in any
%   enumeration); A68 (incoming values are events, handled after the step in
%   progress); A69, A70 (outgoing order per port, host outside); A71 (bad
%   incoming value rejected); A78 (suspension at any andThen); A79, A80, A81
%   (process primitives); A83, A85 (platform parameter, no effect modules);
%   A92 (Bytes flags, boundary values); A96 (main's principal scheme);
%   A86 (the three guarantees and no more); D2, D10, D11.
% - ILLUMINATING: comments on the absence of a rule for failing
%   elaboration, the choice of E over the interface for main, the defaulting
%   of number, self-kill, the stuck message queue when update has no result,
%   input-enabledness, and a program that does not start.
% - TRACTABLE: stated properties: the meaning does not depend on the order
%   chosen among independent modules; the ordering guarantees of A86 follow
%   from the transitions.
% - ECONOMICAL: options merge the silent and platform process steps; one
%   rule ends a process without a message; one rule admits all platform
%   steps.
\chapter{Programs}\label{ch:programs}

This chapter defines what a program is and what it means. A program is a
main module together with the modules it imports. It is elaborated and
evaluated module by module, in a basis that pairs a static basis
(Chapter~\ref{ch:static-modules}) with a dynamic basis
(Chapter~\ref{ch:dynamic-modules}); elaboration finds the type of \kw{main} in
the main module and the ports of the program, and evaluation finds the value
of \kw{main}. That value is a program value (Chapter~\ref{ch:effects}), and the
program runs on the platform (Chapter~\ref{ch:introduction}), which supplies
the primitives and managers beyond those of tasks and ports. % A83
The running of a program is defined as a labelled transition system, whose
configurations hold the current model, the messages awaiting \kw{update}, the
processes, the current subscriptions and the values awaiting delivery on
outgoing ports. The sentences of the chapter infer that a program elaborates
to a type and a port environment, that it evaluates to a value, that a program
value started with flags begins in a configuration, and that a configuration
moves to another, with or without an event. The meaning of a program is the
set of sequences of events at its boundary that it can produce. % D10, D11

\section{Syntax of Programs}\label{sec:programs-syntax}

The phrase class Program of programs is defined as follows:
\begin{center}
\begin{grammar}
\gprod{\program}{\modm\ \opt{\program}}{}
\end{grammar}
\end{center}
A program is thus a non-empty sequence of modules; its last module is its
\emph{main module}. A module $\modm$ \emph{imports} a module $\modid$ if an
import of $\modm$, written or default (Figure~\ref{fig:default-imports}),
names $\modid$. Programs are subject to the following restrictions, in
addition to the restrictions of Section~\ref{sec:syntax-modules-restrictions}
on each module.
\begin{itemize}
\item\citelabel{res:program-modules-distinct} No two modules of a program have the same module identifier, and no
  module of a program has the identifier of a module of the initial basis
  (Appendix~\ref{ch:appendix-static-basis}). % A89
  % Rule r:sm-module's side condition modid notin Dom(ME of B) then holds for
  % every module of a program.
\item\citelabel{res:program-import-order} Each module of a program imports only modules of the initial basis and
  modules that stand before it in the program.
\item\citelabel{res:program-modules-reachable} Every module of a program other than the main module is imported,
  directly or through other modules, by the main module.
  % A16, A30 (amended): every module, the main module included, has a header
  % (Section sec:syntax-modules-restrictions), so no restriction here
  % concerns a module without one.
\item\citelabel{res:program-main} The main module has a top-level value declaration that binds
  \kw{main}. That declaration has the form
  $\kw{main}\ \opt{\kw{:}\ \ty}\ \kw{=}\ \expm$, binding \kw{main} alone,
  and \kw{main} does not occur free in $\expm$, so that \kw{main} does not
  refer to itself, directly or through a delayed occurrence. % A89
\end{itemize}
Thus the modules of a program are the main module and every module it
imports directly or indirectly; the import relation among them has no cycle;
and every module stands after the modules it imports. A program as written is
the finite set of these modules. Any sequence of that set that meets the
restrictions is a program in the sense of the grammar; the main module comes
last in every such sequence, and all such sequences have the same meaning
(Section~\ref{sec:programs-evaluation}).

\section{Elaboration of Programs}\label{sec:programs-elaboration}

Until now, the name Basis and the metavariable $B$ have stood for a static
basis in Chapter~\ref{ch:static-modules} and for a dynamic basis in
Chapter~\ref{ch:dynamic-modules}; those chapters never needed both at once.
In this chapter the class Basis of Chapter~\ref{ch:static-modules} is called
$\pgStaticBasis$, with metavariable $\pgBstat$, and that of
Chapter~\ref{ch:dynamic-modules} is called $\pgDynamicBasis$, with
metavariable $\pgBdyn$; a basis is a pair of the two:
\[
  B\ \text{or}\ (\pgBstat,\pgBdyn)\in\Basis=\pgStaticBasis\times\pgDynamicBasis .
\]
The \emph{initial basis} $B_0$ is the pair of the initial static basis of
Appendix~\ref{ch:appendix-static-basis} and the initial dynamic basis of
Appendix~\ref{ch:appendix-dynamic-basis}. Modification of static bases is
component by component, the set component by union, as for contexts
(Section~\ref{sec:static-core-projection}):
\[
  (T,\ME,\PE)+(T',\ME',\PE')=(T\cup T',\ \ME+\ME',\ \PE+\PE') ;
\]
a dynamic basis is a module environment, modified as a finite map.

The rules of this section and the next allow sentences of the forms
\[
  \begin{array}{c}
  \jstat{\pgBstat}{\program}{(\tau,\PE)},\qquad
  \jdyn{\pgBdyn}{\program}{v},\\[0.5ex]
  \jdg{B}{\program}{(\tau,\PE,v)}
  \end{array}
\]
to be inferred, where $\pgBstat$ is a static basis, $\pgBdyn$ a dynamic basis,
$B$ a basis, $\tau$ a type, $\PE$ a port environment
(Section~\ref{sec:static-modules-objects}) and $v$ a value. We write
$\vdash_{\mathrm{STAT}}$ for the elaboration of programs, $\vdash_{\mathrm{DYN}}$
for their evaluation, and $\vdash$ alone for the sentence that combines the
two. The first sentence is pronounced ``$\program$ elaborates to $\tau$ and
$\PE$ in $\pgBstat$'': $\tau$ is the type of \kw{main} in the main module, and
$\PE$ holds the ports of the program. The second is pronounced ``$\program$
evaluates to $v$ in $\pgBdyn$'': $v$ is the value of \kw{main}. The third is
pronounced ``$\program$ elaborates to $\tau$ and $\PE$, and evaluates to $v$,
in $B$''. Premises also use the sentences about modules of
Chapters~\ref{ch:static-modules} and~\ref{ch:dynamic-modules}, in their forms
there; other hypotheses are side conditions.

The type of \kw{main} is the instance of its type scheme that takes $\tFloat$
for each bound type variable of class $\knum$. It is the type of the default
copy of the declaration of \kw{main}, the copy that number specialisation
makes of every top-level declaration, taking $\tFloat$ for each of its variables
of class $\knum$ (Section~\ref{sec:dynamic-core-numbers}). % A27, A87
For a type scheme $\sigma=\forall\tyrowvar^{(k)}.\tau$, let
$\alpha_1,\ldots,\alpha_l$ be those of the bound variables
$\tyrowvar_1,\ldots,\tyrowvar_k$ that are type variables of class $\knum$; then
\[
  \pgDflt\sigma=\tau\{\tFloat,\ldots,\tFloat/\alpha_1,\ldots,\alpha_l\} .
\]
The substitution respects classes, since $\tFloat\in\Adm{\knum}$. In the rules
below, $\sbtProgram$ is the type name of \kw{Platform.Program}, of arity 3,
and $\sbtNode$ that of \kw{VirtualDom.Node}, of arity 1, of which \kw{Html} is
an alias (Appendix~\ref{ch:appendix-static-basis}); $\effMgrPlat$ is the set
of the names of the managers of the platform
(Section~\ref{sec:effects-simple}). In each rule, the premise
$\jdg{\pgBstat}{\modm}{\pgBstat',E}$ is inferred by a principal derivation for
the module (Section~\ref{sec:static-modules-principal}), so that $E$, and the
scheme $\sigma$ of \kw{main} in it, are principal. % A96

\rulegroup{Elaboration of Programs}{\jstat{\pgBstat}{\program}{(\tau,\PE)}}

\defnrule[r:p-stat-module]
  {\jdg{\pgBstat}{\modm}{\pgBstat',E}\\
   \Dom(\PE\of\pgBstat')\cap\big(\Dom(\PE\of\pgBstat)\cup\effMgrPlat\big)=\emptyset\\
   \jstat{\pgBstat+\pgBstat'}{\program}{(\tau,\PE)}}
  {\jstat{\pgBstat}{\modm\ \program}{(\tau,\PE)}}

\defnrule[r:p-stat-main-program]
  {\jdg{\pgBstat}{\modm}{\pgBstat',E}\\
   \Dom(\PE\of\pgBstat')\cap\big(\Dom(\PE\of\pgBstat)\cup\effMgrPlat\big)=\emptyset\\
   E(\kw{main})=(\sigma,\isv)\premsep
   \sigma=\forall\tyrowvar^{(k)}.(\tau_1,\tau_2,\tau_3)\sbtProgram\premsep
   \Port(\tau_1)}
  {\jstat{\pgBstat}{\modm}{\big(\pgDflt\sigma,\ \PE\of(\pgBstat+\pgBstat')\big)}}

\defnrule[r:p-stat-main-view]
  {\jdg{\pgBstat}{\modm}{\pgBstat',E}\\
   \Dom(\PE\of\pgBstat')\cap\big(\Dom(\PE\of\pgBstat)\cup\effMgrPlat\big)=\emptyset\\
   E(\kw{main})=(\sigma,\isv)\premsep
   \pgDflt\sigma=\tau\,\sbtNode}
  {\jstat{\pgBstat}{\modm}{\big(\tau\,\sbtNode,\ \PE\of(\pgBstat+\pgBstat')\big)}}

\begin{rulecomments}
\rcomment{r:p-stat-module} Each module is elaborated in the static basis that
the modules before it have built, so it finds there the interface of every
module it imports, and the type names it generates are new
(rule~\ruleref{r:sm-module}). The side condition makes port names global to
the program: no two ports of a program share a name, even in different
modules, and no port has the name of a manager of the platform; a program
with such a clash does not elaborate. % A84, D2
\rcomment{r:p-stat-main-program} Rules~\ruleref{r:p-stat-main-program}
and~\ruleref{r:p-stat-main-view} cover the one production, since \kw{main}
may have either of two forms of type. Only the \kw{main} of the main module
is constrained; a \kw{main} in another module is an ordinary value. % A32
The scheme of \kw{main} is taken from $E$, the environment of all the
declarations of the module, not from its interface, so \kw{main} need not be
exposed. % A82
A bound variable of class $\knum$ in the scheme of \kw{main} is $\tFloat$ in
its type: with \kw{init = \textbackslash() -> (0, Cmd.none)} and no other
constraint on the model, the model type is $\tFloat$, and the value of
\kw{main} is that of its default copy, at $\tFloat$
(rule~\ruleref{r:p-dyn-main}). % A27, A87
The flags type $\tau_1$ is a port type, and is taken from the scheme before
the default, so it contains no type variable. Thus
\kw{main : Program (Int -> Int) Model Msg} and
\kw{main : Program flags Model Msg}, with \kw{flags} a type variable, do not
elaborate, and nor does an unannotated \kw{main} whose \kw{init} is
\kw{\textbackslash n -> (n + 1, Cmd.none)}, whose flags type is a variable of
class $\knum$. % A96
The unit type is a port type, and as flags it is
treated specially (Section~\ref{sec:programs-execution}). Since \kw{Bytes}
is a port type, flags may be, or contain, values of type \kw{Bytes}. % A92
The model type $\tau_2$ and the message type $\tau_3$ are unconstrained. Some
implementations admit further types of flags, which are not port types
(Appendix~\ref{ch:appendix-implementation}).
% A33, A77
\rcomment{r:p-stat-main-view} A \kw{main} of type \kw{Html msg} or
\kw{Svg msg} has a type $\tau\,\sbtNode$, since both are aliases. No function
type is of either form, so a \kw{main} with parameters does not elaborate. % A82
\end{rulecomments}

Since the derivations for the modules of a program are principal, a flags
type that the program leaves undetermined is a type variable in the scheme of
\kw{main}, which is not a port type: an unannotated
\kw{main = Platform.worker \{ init = \textbackslash\_ -> (0, Cmd.none), \ldots\ \}},
whose \kw{init} ignores its flags, does not elaborate, although a derivation
that took the flags to be $\tunit$ would satisfy the rule. % A96, A72, A82

There is no rule for a program one of whose modules does not elaborate. Such
a program has no meaning, and no part of it is evaluated
(Section~\ref{sec:programs-evaluation}).

\section{Evaluation of Programs}\label{sec:programs-evaluation}

The modules of a program are evaluated in the order in which they are
elaborated, each in the dynamic basis that the modules before it have built.
The syntax evaluated is the reduced syntax of
Section~\ref{sec:dynamic-modules-reduced}, which records the direction and
payload type that elaboration finds for each port, after number
specialisation (Section~\ref{sec:dynamic-core-numbers}), which makes the
default copy of every top-level declaration besides the copies its uses
need; when the declaration quantifies no variable of class $\knum$, the
default copy is the declaration itself. In rule~\ruleref{r:p-dyn-main},
$\VE(\kw{main})$ is the value bound by the default copy of the declaration of
\kw{main}, whose type is $\pgDflt\sigma$ for the scheme $\sigma$ of \kw{main},
the type found by rule~\ruleref{r:p-stat-main-program}
or~\ruleref{r:p-stat-main-view}. % A27, A87

\rulegroup{Evaluation of Programs}{\jdyn{\pgBdyn}{\program}{v}}

\defnrule[r:p-dyn-module]
  {\jdg{\pgBdyn}{\modm}{\pgBdyn',\VE}\premsep
   \jdyn{\pgBdyn+\pgBdyn'}{\program}{v}}
  {\jdyn{\pgBdyn}{\modm\ \program}{v}}

\defnrule[r:p-dyn-main]
  {\jdg{\pgBdyn}{\modm}{\pgBdyn',\VE}\premsep v=\VE(\kw{main})}
  {\jdyn{\pgBdyn}{\modm}{v}}

\begin{rulecomments}
\rcomment{r:p-dyn-module} A module is evaluated once, after every module it
imports, and its top-level values are evaluated then, whether or not they are
used: every top-level declaration has its default copy, even when nothing
uses it. A \kw{main} in a module other than the main module is evaluated in the
same way, once, as an ordinary value. % A59, A32, A87
\rcomment{r:p-dyn-main} The value of \kw{main} is taken from $\VE$, the value
environment of all the top-level declarations of the module
(rule~\ruleref{r:dm-module}), not from the basis $\pgBdyn'$, which holds only
what the module exposes; so \kw{main} need not be exposed. No module imports
the main module, and $\pgBdyn'$ is not used. % A82
If a top-level value of any module has no value (it reaches
\kw{Debug.todo}, or forces a delayed reference to itself), no sentence about
the program can be inferred, and the program has no meaning. % A25, A62
\end{rulecomments}

\rulegroup{Programs}{\jdg{B}{\program}{(\tau,\PE,v)}}

\defnrule[r:p-program]
  {\jstat{\pgBstat\of B}{\program}{(\tau,\PE)}\premsep
   \jdyn{\pgBdyn\of B}{\program}{v}}
  {\jdg{B}{\program}{(\tau,\PE,v)}}

\begin{rulecomments}
\rcomment{r:p-program} A program is elaborated as a whole before any part of
it is evaluated: if one module fails to elaborate, the static premise fails,
and nothing is evaluated, not even the modules before it. Elaboration may
choose the new type names of a module freely (rule~\ruleref{r:sm-module}), but
neither $\PE$, whose types are built from type names of the initial basis
alone, nor $v$ depends on the choice.
\end{rulecomments}

A program $\program$ \emph{has a meaning} when
$\jdg{B_0}{\program}{(\tau,\PE,v)}$ for some $\tau$, $\PE$ and $v$. Since a
program is a set of modules, the following property, which may easily be
demonstrated, is needed for this to make sense.
\begin{property}
Let $\program$ and $\program'$ be sequences of the same modules, both meeting
the restrictions of Section~\ref{sec:programs-syntax}. Then the sentence
$\jdg{B_0}{\program}{(\tau,\PE,v)}$ can be inferred if and only if the
sentence $\jdg{B_0}{\program'}{(\tau',\PE,v)}$ can, where $\tau'$ is $\tau$
with the type names that the modules generate renamed.
\end{property}

\section{Execution}\label{sec:programs-execution}

Let $\jdg{B_0}{\program}{(\tau,\PE,v)}$. When $\tau$ is $\tau'\,\sbtNode$, the
program is a view, whose meaning is given in
Section~\ref{sec:programs-views}. Otherwise
$\tau=(\tau_1,\tau_2,\tau_3)\sbtProgram$, where $\tau_1$ is the \emph{flags
type}, and $v$ is a program value $\effprog\,r$
(Section~\ref{sec:effects-commands}). This section defines how a program
value runs. Throughout, $\PE$ is the port environment of the program, and the
record $r$ is fixed: the rules apply its fields \kw{init}, \kw{update} and
\kw{subscriptions}, written $r(\kw{init})$ and so on.

A program runs on a platform. Besides the primitives $\effPrimPlat$, their
outcomes and the managers $\effMgrPlat$ (Chapter~\ref{ch:effects}), the
platform provides a set $\pgEventPlat$ of \emph{platform events}, and its
\emph{platform steps}, defined in Section~\ref{sec:programs-transitions}. All
of these are parameters of the platform; the definition fixes the manager of
tasks, the ports, and the program loop. % A83, A85

\subsection{Configurations and Events}\label{sec:programs-configurations}

For a set $A$, $A^{*}=\bigcup_{n\geq 0}A^n$ is the set of finite sequences of
members of $A$. The empty sequence is $\pgnil$, and $x\cdot y$ is the
concatenation of the sequences $x$ and $y$; a member of $A$ is also a
sequence of length one, and $x_1\cdots x_n$ is $x_1\cdot\ldots\cdot x_n$.
A finite multiset of members of $A$ is a finite map from $A$ to the positive
integers, giving the multiplicity of each member; multisets are written as in
Section~\ref{sec:effects-compound}. The objects of execution are shown in
Figure~\ref{fig:programs-objects}.

\begin{figure}[htbp]
\centering
\begin{objtable}
\objdef{s\ \text{or}\ d}{\pgTime}{\{x\in\mathbb{R}\mid x\geq 0\}}
\objdef{K}{\Config}{\pgTime\times\Val\times\Val^{*}\times\pgProcMap}
& & & \phantom{{}={}}\times\SubVal\times\pgOutQueue\times\pgPending & \\
\objdef{P}{\pgProcMap}{\ProcId\finmap\pgProc}
\objdef{}{\pgProc}{\big(\TaskVal\times\pgTime\times(\CmdVal\cup\{\pgspawned\})\big)}
& & & \phantom{{}={}}\cup\{\pgended\} & \\
\objdef{Q}{\pgOutQueue}{\cls{Var}\finmap\BVal^{*}}
\objdef{H}{\pgPending}{\CmdVal\finmap\{1,2,\ldots\}}
\objdef{\ell}{\Event}{(\{\pgEvIn\}\times\cls{Var}\times\BVal)}
& & & \phantom{{}={}}\cup(\{\pgEvOut\}\times\cls{Var}\times\BVal)\cup\effAction & \\
& & & \phantom{{}={}}\cup(\{\pgwait\}\times\pgTime)\cup\pgEventPlat & \\
\end{objtable}
\caption{Semantic Objects of Execution}
\label{fig:programs-objects}
\end{figure}

A configuration $K=(s,v,M,P,\sbv,Q,H)$ has the following components.
\begin{itemize}
\item $s$ is the present time, in milliseconds since the program started.
\item $v$ is the \emph{model}, the value that the latest \kw{update}, or
  \kw{init}, produced.
\item $M$ is the \emph{message queue}, the messages awaiting \kw{update}.
\item $P$ holds every process started so far. If
  $P(\pid)=(\tk,s',\pgorig)$, the process named by $\pid$ is \emph{live}, is
  reducing $\tk$, and took its latest step (or started) at time $s'$; its
  \emph{origin} $\pgorig$ is the leaf command whose task it performs, or
  $\pgspawned$ if it was started by $\effspawn$. If $P(\pid)=\pgended$, the
  process has ended or been killed.
\item $\sbv$ is the current subscription value, the value of
  \kw{subscriptions} applied to the model.
\item $Q(\effp)$ is the sequence of boundary values awaiting delivery on the
  outgoing port $\effp$.
\item $H$ holds the leaf commands handed to the managers of the platform and
  not yet acted on.
\end{itemize}
The event $\pgEvIn(\effp,j)$ is the host's sending of $j$ on the incoming port
$\effp$; $\pgEvOut(\effp,j)$ is the program's delivery of $j$ on the outgoing
port $\effp$; an action $\effa$ (Section~\ref{sec:effects-tasks}) is a
primitive performed by a process; $\pgwait\,d$, with $d>0$, is the passing of
$d$ milliseconds; and a platform event is whatever the platform makes
observable of its managers.

For a manager $\effm$ and a value $v$, the leaf commands (and leaf
subscriptions) of $\effm$ with argument $v$ are
\[
  \pgLeaf(\effm,v)=\{\effmap(v_n,\cdots\effmap(v_1,\effleaf(\effm,v))\cdots)
    \mid n\geq 0\} .
\]
The outgoing queues at the start of a program are
\[
  \pgQinit=\{\effp\mapsto\pgnil\,;\,\PE(\effp)=(\smdirout,\tau)\ \text{for some}\ \tau\} .
\]
For a type $\tau$ with $\Port(\tau)$ and a boundary value $j$, the \emph{flags
value} of $j$ at $\tau$ is
\[
  \pgFlags{\tau}(j)=\begin{cases}
    () & \text{if}\ \tau=\tunit,\\
    \Dec{\tau}(j) & \text{otherwise},
  \end{cases}
\]
which is undefined when $\Dec{\tau}(j)$ is. % A77

\subsection{Dispatch and Start}\label{sec:programs-dispatch}

The rules of this section and the next allow sentences of the forms
\[
  \jdg{s,P,Q,H}{L}{(P',Q',H')},\qquad
  \jdg{\PE,\tau_1,j}{\effprog\,r}{K},\qquad
  K\step{\opt{\ell}}K'
\]
to be inferred. In the first, $L$ is a sequence of leaf commands, and the
sentence is pronounced ``dispatching $L$ at time $s$ turns $P$, $Q$, $H$ into
$P'$, $Q'$, $H'$''. The second is pronounced ``$\effprog\,r$, started with
the flags $j$, begins in $K$''. The third is pronounced ``$K$ moves to $K'$'',
and, when the event $\ell$ is present, ``with $\ell$''; a transition without
an event is \emph{silent}. Premises also use the sentences of
Chapter~\ref{ch:effects} and the sentence $\japp{v_1}{v_2}{v}$ of the
application of values (Section~\ref{sec:dynamic-core-apply}); other
hypotheses are side conditions. The convention for options is the same as for
the Core static semantics. There are no states and no exception packets: the
configuration is the whole of what changes, and every rule is given in full.
% D2

A command is \emph{dispatched} by dispatching a sequence that holds its leaf
commands, each as many times as $\effFlat$ gives it, in any order. Which
order is no part of the meaning, by the effect order convention of
Section~\ref{sec:effects-commands}. % A66, A67

\rulegroup{Dispatch}{\jdg{s,P,Q,H}{L}{(P',Q',H')}}

\defnrule[r:p-disp-empty]
  {}
  {\jdg{s,P,Q,H}{\pgnil}{(P,Q,H)}}

\defnrule[r:p-disp-task]
  {\effc\in\pgLeaf(\effmtask,\tk)\premsep\pid\notin\Dom P\\
   \jdg{s,\ P+\{\pid\mapsto(\tk,s,\effc)\},\ Q,H}{L}{(P',Q',H')}}
  {\jdg{s,P,Q,H}{\effc\cdot L}{(P',Q',H')}}

\defnrule[r:p-disp-out]
  {\effc\in\pgLeaf(\effmout\,\effp,j)\\
   \jdg{s,P,\ Q+\{\effp\mapsto Q(\effp)\cdot j\},\ H}{L}{(P',Q',H')}}
  {\jdg{s,P,Q,H}{\effc\cdot L}{(P',Q',H')}}

\defnrule[r:p-disp-plat]
  {\effc\in\pgLeaf(\effm,v)\premsep\effm\in\effMgrPlat\\
   \jdg{s,P,Q,\ H\uplus\effmset{\effc}}{L}{(P',Q',H')}}
  {\jdg{s,P,Q,H}{\effc\cdot L}{(P',Q',H')}}

\begin{rulecomments}
\rcomment{r:p-disp-task} Each task leaf starts a process of its own, with an
identifier that names no process started before, live or ended; a leaf that
occurs twice in a command starts two processes, and performs its task twice.
% D10
The process keeps the leaf as its origin, since the message it gives, when
its task succeeds, comes from the leaf (rule~\ruleref{r:p-proc-message}).
\rcomment{r:p-disp-out} The boundary value was computed when the command was made
(Section~\ref{sec:effects-ports}); dispatch only appends it to the queue of
its port. Since $Q(\effp)$ is defined only for the outgoing ports of the
program, $\effp$ is one of them. The values of one command for one port
join the queue in the order of $L$, which is unspecified: with
\kw{Cmd.batch [send 1, send 2]}, either value may be delivered first. % A66
\rcomment{r:p-disp-plat} What the managers of the platform do with the leaves
handed to them is given by the platform steps
(Section~\ref{sec:programs-transitions}). No rule dispatches a leaf of any
other manager: a leaf of an incoming port is a subscription, never a
command, and there are no effect modules in Standard Elm. % A83, A85
\end{rulecomments}

\rulegroup{Starting a Program}{\jdg{\PE,\tau_1,j}{\effprog\,r}{K}}

\defnrule[r:p-init]
  {\pgFlags{\tau_1}(j)=v_1\premsep\japp{r(\kw{init})}{v_1}{(v,\cm)}\\
   \effFlat(\cm)=\effmset{\effc_1,\ldots,\effc_n}\premsep
   \jdg{0,\ \emptymap,\ \pgQinit,\ \emptymap}{\effc_1\cdots\effc_n}{(P,Q,H)}\\
   \japp{r(\kw{subscriptions})}{v}{\sbv}}
  {\jdg{\PE,\tau_1,j}{\effprog\,r}{(0,v,\pgnil,P,\sbv,Q,H)}}

\begin{rulecomments}
\rcomment{r:p-init} The host starts a program with a boundary value $j$ as
its flags; when the host gives no flags, $j$ is $\effnull$. % A72, A92
With flags of the unit type every boundary value is accepted, whereas a port of
the unit type accepts only $\effnull$. % A77
When $\pgFlags{\tau_1}(j)$ is undefined, as for $j=1.5$ with flags type
\kw{Int}, no rule applies, there is no configuration in which the program
begins, and the program does not start; how an implementation reports this
is outside the definition (Appendix~\ref{ch:appendix-implementation}). % D2
The commands of \kw{init} are dispatched, and the subscriptions computed from
the initial model, before any message is handled. The empty map
$\emptymap$ serves here as the empty process map and the empty multiset.
\end{rulecomments}

\subsection{Transitions}\label{sec:programs-transitions}

A \emph{platform step} $K\pgplatstep{\opt{\ell}}K'$ is given by the platform.
Its event $\ell$, if present, is a platform event. A platform step changes
nothing but the components $M$ and $H$ of $K$: it removes leaf commands from
$H$, and it appends to $M$ messages that, as the platform defines, are given
by the leaf commands it removes or by the members of $\effFlat(\sbv)$ whose
manager is in $\effMgrPlat$ (Section~\ref{sec:effects-messages}). Whether a
platform step exists may depend on the present time $s$. % A83

\rulegroup{Transitions}{K\step{\opt{\ell}}K'}

\defnrule[r:p-update]
  {\japp{r(\kw{update})}{v_1}{v_2}\premsep\japp{v_2}{v}{(v',\cm)}\\
   \effFlat(\cm)=\effmset{\effc_1,\ldots,\effc_n}\premsep
   \jdg{s,P,Q,H}{\effc_1\cdots\effc_n}{(P',Q',H')}\\
   \japp{r(\kw{subscriptions})}{v'}{\sbv'}}
  {(s,v,v_1\cdot M,P,\sbv,Q,H)\step{}(s,v',M,P',\sbv',Q',H')}

\defnrule[r:p-proc-step]
  {P(\pid)=(\tk,s',\pgorig)\premsep\tk\effred{\opt{\effa}}\tk'\\
   \opt{\effa=\effdone{\pi(v_1\cdots v_k)}{\tk''}\premsep\pi\in\effPrimPlat}}
  {(s,v,M,P,\sbv,Q,H)\step{\opt{\effa}}
   (s,v,M,P+\{\pid\mapsto(\tk',s,\pgorig)\},\sbv,Q,H)}

\defnrule[r:p-proc-sleep]
  {P(\pid)=(\tk,s',\pgorig)\premsep\tk\effred{\effa}\tk'\premsep
   \effa=\effdone{\effsleep(v')}{\tk''}\\
   \text{if}\ v'\in\effFltFin\ \text{then}\ s\geq s'+\max(0,v')}
  {(s,v,M,P,\sbv,Q,H)\step{\effa}
   (s,v,M,P+\{\pid\mapsto(\tk',s,\pgorig)\},\sbv,Q,H)}

\defnrule[r:p-proc-spawn]
  {P(\pid)=(\tk,s',\pgorig)\premsep\tk\effred{\effa}\tk'\premsep
   \effa=\effdone{\effspawn(\tk_0)}{\effsucceed\,\pid'}\\
   \pid'\notin\Dom P\premsep
   P'=P+\{\pid\mapsto(\tk',s,\pgorig),\ \pid'\mapsto(\tk_0,s,\pgspawned)\}}
  {(s,v,M,P,\sbv,Q,H)\step{\effa}(s,v,M,P',\sbv,Q,H)}

\defnrule[r:p-proc-kill]
  {P(\pid)=(\tk,s',\pgorig)\premsep\tk\effred{\effa}\tk'\premsep
   \effa=\effdone{\effkill(\pid')}{\tk''}\\
   P'=(P+\{\pid'\mapsto\pgended\})+\{\pid\mapsto(\tk',s,\pgorig)\}}
  {(s,v,M,P,\sbv,Q,H)\step{\effa}(s,v,M,P',\sbv,Q,H)}

\defnrule[r:p-proc-message]
  {P(\pid)=(\effsucceed\,v',s',\effc)\premsep\effrecv{\effc}{v'}{v''}}
  {(s,v,M,P,\sbv,Q,H)\step{}(s,v,M\cdot v'',P+\{\pid\mapsto\pgended\},\sbv,Q,H)}

\defnrule[r:p-proc-end]
  {P(\pid)=(\tk,s',\pgorig)\\
   \tk=\efffail\,v'\ \text{or}\ (\tk=\effsucceed\,v'\ \text{and}\ \pgorig=\pgspawned)}
  {(s,v,M,P,\sbv,Q,H)\step{}(s,v,M,P+\{\pid\mapsto\pgended\},\sbv,Q,H)}

\defnrule[r:p-in-accept]
  {\PE(\effp)=(\smdirin,\tau)\premsep\Dec{\tau}(j)=v'\premsep
   \effFlat(\sbv)=\effmset{\effc_1,\ldots,\effc_n}\\
   \text{for}\ 1\leq i\leq n:\ \ \effrecv{\effc_i}{(\effp,j)}{M_i},\\
   \text{or}\ M_i=\pgnil\ \text{and}\
   \effrecv{\effc_i}{(\effp,j)}{v''}\ \text{for no}\ v''}
  {(s,v,M,P,\sbv,Q,H)\step{\pgEvIn(\effp,j)}
   (s,v,M\cdot M_1\cdots M_n,P,\sbv,Q,H)}

\defnrule[r:p-in-reject]
  {\PE(\effp)=(\smdirin,\tau)\premsep\Dec{\tau}(j)\ \text{is undefined}}
  {K\step{\pgEvIn(\effp,j)}K}

\defnrule[r:p-out]
  {Q(\effp)=j\cdot J}
  {(s,v,M,P,\sbv,Q,H)\step{\pgEvOut(\effp,j)}(s,v,M,P,\sbv,Q+\{\effp\mapsto J\},H)}

\defnrule[r:p-wait]
  {d>0}
  {(s,v,M,P,\sbv,Q,H)\step{\pgwait\,d}(s+d,v,M,P,\sbv,Q,H)}

\defnrule[r:p-platform]
  {K\pgplatstep{\opt{\ell}}K'}
  {K\step{\opt{\ell}}K'}

\begin{rulecomments}
\rcomment{r:p-update} This is the only rule that takes a message from the
queue or changes the model. The message at the head of the queue is applied,
with the model, to \kw{update}; the commands of the result are dispatched,
and the subscriptions recomputed from the new model, whether or not it has
changed, all in the one step. If the application of \kw{update} has no
result (it reaches \kw{Debug.todo}), the rule never applies again, and no
further message is handled, though processes, ports and the platform go on
taking steps. % A62
\rcomment{r:p-proc-step} A process takes one reduction of its task at a time,
and the steps of different processes, and every other transition, may come
between any two of them. A process may thus be suspended at any
$\effandThen$; a program may rely on nothing more. % A78
Rules~\ruleref{r:p-proc-step} to~\ruleref{r:p-proc-kill} record the present
time as the time of the process's latest step. An action of a primitive of
the platform acts on the world as the platform defines.
\rcomment{r:p-proc-sleep} The process reached $\effprim(\effsleep,v')$ by its
step at time $s'$, so the action takes place no sooner than $\max(0,v')$
milliseconds later; for an infinite or NaN delay, no condition is imposed.
The definition sets no time by which the action must take place. % A81
\rcomment{r:p-proc-spawn} The new process names a process never started
before, and its origin is $\pgspawned$: its outcome goes to no other process,
and it gives no message (rule~\ruleref{r:p-proc-end}). % A79
\rcomment{r:p-proc-kill} The killed process takes no further step, and if it
was performing the task of a leaf it gives no message. Killing a process
that has ended changes nothing. When a process kills itself, the second
modification overrides the first, so it is not ended and goes on with
$\tk'$. % A80
\rcomment{r:p-proc-message} A process whose origin is a leaf command ends,
when its task succeeds, by adding the message of its leaf
(rule~\ruleref{r:e-msg-task}) to the end of the queue. When the functions of
the $\effmap$s around the leaf give no result, the process never ends.
\rcomment{r:p-proc-end} A process whose task fails ends, whatever its origin,
and nothing else happens: there is no error, and no message. % A79, D2
\rcomment{r:p-in-accept} An incoming value is an event of its own, so it is
never handled in the middle of another step; its messages join the end of the
queue, and are handled by later updates. % A68
Each leaf subscription on $\effp$ gives one message, through the functions of
the $\effmap$s around it; a leaf on another port, or of another manager,
gives none, and so does a leaf whose functions give no result. The order of
the messages of one value follows the enumeration of $\effFlat(\sbv)$, and
is unspecified. A value accepted while no subscription is current gives no
message. % A66
\rcomment{r:p-in-reject} A boundary value that does not decode at the type of the
port is rejected: the program never sees it, and is unchanged. % A71, D2
Thus with \kw{port recv : (String -> msg) -> Sub msg}, the value 1337 is
rejected. For every incoming port and every boundary value one of
rules~\ruleref{r:p-in-accept} and~\ruleref{r:p-in-reject} applies, so the host
may send at any time. A name that is not an incoming port of the program has
no transition: sending on it is outside the definition.
\rcomment{r:p-out} The values of each port are delivered in the order in
which they joined its queue, so a value sent by one update is delivered
before a value that a later update sends on the same port. Delivery happens
after the step that queued the value, at a time the definition does not fix,
and in no fixed order across ports. The event records only that the program
has delivered $j$; whether the host has a subscriber, and what it does with
$j$, are outside the definition. % A69, A70
\rcomment{r:p-wait} Time passes only in this rule, and every other transition
takes no time. Since this rule applies to every configuration, the definition
requires no step to happen by any time: only the lower bound of
rule~\ruleref{r:p-proc-sleep}, and what the platform imposes on its own
steps, relate the steps of a program to time.
\rcomment{r:p-platform} The managers of the platform, such as those of
\kw{Time} or \kw{Http}, act through these steps, within the limits stated
above. A platform with no managers has no platform steps. % A83
\end{rulecomments}

The rules give the following guarantees, and no others. % A86
\begin{property}
Messages are handled one at a time, by rule~\ruleref{r:p-update}, in the
order in which they joined the message queue. Each update applies
\kw{update} to the model that the previous update produced, or to the model
of \kw{init} for the first. The commands of an update are dispatched in the
step of that update, so the effects of an earlier update are dispatched
before those of a later one. Apart from these, and from the order of the
values of each outgoing port, the steps of distinct processes, the arrival
of incoming values and platform steps, the delivery of outgoing values, and
the update steps interleave in every way the rules allow.
\end{property}

\section{Observable Behaviour}\label{sec:programs-behaviour}

Let $\jdg{B_0}{\program}{(\tau,\PE,\effprog\,r)}$ with
$\tau=(\tau_1,\tau_2,\tau_3)\sbtProgram$, and let $j$ be a boundary value. An
\emph{execution} of $\program$ with flags $j$ is a finite sequence of
transitions
\[
  K_0\step{\opt{\ell_1}}K_1\step{\opt{\ell_2}}\cdots\step{\opt{\ell_n}}K_n
  \qquad(n\geq 0),
\]
where $\jdg{\PE,\tau_1,j}{\effprog\,r}{K_0}$. The \emph{observable} events
are those of the forms $\pgEvIn(\effp,j')$, $\pgEvOut(\effp,j')$ and
$\pgwait\,d$, the platform events, and the actions
$\effdone{\pi(v_1\cdots v_k)}{\tk}$ with $\pi\in\effPrimPlat$. The
\emph{trace} of an execution is the sequence of its observable events, in
order; silent transitions and the actions of $\effsleep$, $\effspawn$ and
$\effkill$, whose whole effect lies within the configuration, are omitted.
The set of traces
\[
  \pgTr(\program,j)
\]
is the set of the traces of the executions of $\program$ with flags $j$, and
the \emph{meaning} of $\program$ is the function that maps each $j$ to
$\pgTr(\program,j)$. When $\pgTr(\program,j)$ is empty, the program does not
start with flags $j$; otherwise it holds the empty trace and every prefix of
each of its traces.

An implementation \emph{conforms} to the definition if, for every program
that has a meaning and every boundary value $j$, every trace that the
implementation can produce when it runs the program with flags $j$ is a
member of $\pgTr(\program,j)$, and the implementation does not start the
program when $\pgTr(\program,j)$ is empty. Two conforming implementations may
produce different traces of one program, as where the order of the commands
of a batch, or the interleaving of processes, is left open, and both conform.
Conformance concerns what is produced, not that anything is: the definition
does not require a program to make progress. % A86

\section{Views and Other Programs}\label{sec:programs-views}

Let $\jdg{B_0}{\program}{(\tau'\,\sbtNode,\PE,v)}$. The program is a
\emph{view}: its meaning is the value $v$ of type \kw{Html}, which the host
renders. % A82
It has no configuration and no transitions: it handles no message, and has no
traffic on its ports, whatever ports its modules declare. How the host
renders $v$ is outside the definition; an implementation that only evaluates
the program conforms. Every top-level value of every module is evaluated, so
a \kw{Debug.log} in the definition of \kw{main} is evaluated once.

A platform may provide functions other than \kw{Platform.worker} that make
program values, such as the functions \kw{sandbox}, \kw{element},
\kw{document} and \kw{application} of the module \kw{Browser}. They are not part of the
initial basis, and their records have further fields, such as \kw{view}
(Section~\ref{sec:effects-commands}). A program whose \kw{main} is made by one
of them runs as in Section~\ref{sec:programs-execution}, except that the
platform defines how \kw{init} is applied (for \kw{application}, to
the flags, a URL and a navigation key), and that among its platform steps are
steps that render the value of \kw{view} applied to the present model, as a
platform event, after each update, and steps that append to the message
queue the messages that the host sends from the rendered view. % A83
A platform that does not provide these functions, or does not run the
programs they make, is recorded in Appendix~\ref{ch:appendix-implementation}.
